\documentclass[a4paper,twocolumn]{article}

\RequirePackage[croatian]{babel}
\makeatletter
\renewcommand\thesection{\@arabic\c@section.}
\renewcommand\thesubsection{\thesection\@arabic\c@subsection.}
\renewcommand\thesubsubsection{\thesubsection\@arabic\c@subsubsection.}
\renewcommand\theequation{\@arabic\c@equation}
\renewcommand\thefigure{\@arabic\c@figure.}
\renewcommand\thetable{\@arabic\c@table.}
\makeatother

\RequirePackage{graphicx}
\RequirePackage{hyperref}
\RequirePackage{amssymb}
\RequirePackage{amsmath}
\RequirePackage{booktabs}

\title{Raspoznavanje znakova znakovnog jezika primjenom dubokih neuronskih mreža}
\author{Lucija Čorak, Mateo Jakšić, Petar Kuruc, Martina Šola, Klara Zagajski}
\date{}

\begin{document}

\maketitle

\section*{Sažetak}

U ovom radu istražuju se metode klasifikacije znakova znakovnog jezika korištenjem dubokih neuronskih mreža. Razmatrani su vlastiti CNN model te tri unaprijed istrenirana modela: MobileNetV2, EfficientNetB0 i ResNet50. Usporedba se provodi na temelju metrika: točnost, preciznost, odziv, F1-mjera, broj parametara i vrijeme treniranja.

\section{Uvod}

Znakovni jezik predstavlja osnovno sredstvo komunikacije za gluhe i nagluhe osobe. Automatizirano prepoznavanje znakova znakovnog jezika omogućuje razvoj inkluzivnih tehnologija i alata za lakšu komunikaciju. U ovom radu uspoređuju se četiri duboka modela trenirana na istom skupu podataka, s ciljem analize učinkovitosti pojedinih arhitektura.

\section{Korišteni skup podataka}

Korišten je javno dostupan skup podataka znakova američkog znakovnog jezika (ASL), koji se sastoji od slika koje prikazuju ruke u različitim pozicijama koje odgovaraju slovima engleske abecede. Skup je podijeljen na trening, validacijski i testni dio.

Skup podataka dostupan je na platformi Kaggle: 
\href{https://www.kaggle.com/datasets/datamunge/sign-language-mnist}{Sign Language MNIST} \cite{asl_dataset}.

Radi se o slikama dimenzije $28 \times 28$ piksela u sivim tonovima, a svaka slika prikazuje jednu gestu koja odgovara slovu engleske abecede (osim slova \textit{J} i \textit{Z}, koja uključuju pokret i nisu zastupljena u ovom skupu).

Podaci su raspodijeljeni na sljedeći način:
\begin{itemize}
    \item Trening skup: X uzoraka
    \item Validacijski skup: Y uzoraka
    \item Testni skup: Z uzoraka
\end{itemize}


\section{Obrada slike}

Prije treniranja modela provedena je obrada ulaznih slika kako bi se poboljšala kvaliteta podataka i omogućila robusnost modela.

\subsection{Promjena dimenzija i skaliranje}

Sve slike su skalirane na istu dimenziju $X \times X$ piksela kako bi se zadovoljili ulazni zahtjevi mreža. Pikselske vrijednosti su normalizirane na raspon $[0,1]$.

\subsection{Augmentacija podataka}

Kako bi se smanjilo pretreniranost i povećala generalizacija modela, korištena je augmentacija slika. Primijenjene su sljedeće transformacije:
\begin{itemize}
    \item Horizontalna i vertikalna translacija
    \item Rotacija (do $X^\circ$)
    \item Zrcaljenje (horizontalno)
    \item Zoomiranje (unutar raspona $[Y_1, Y_2]$)
\end{itemize}

\subsection{Podjela skupa}

Skup podataka podijeljen je na tri dijela:
\begin{itemize}
    \item Trening skup: X\% uzoraka
    \item Validacijski skup: Y\% uzoraka
    \item Testni skup: Z\% uzoraka
\end{itemize}

\subsection{Standardizacija prema modelu}

Za pretrenirane modele (MobileNetV2, EfficientNetB0, ResNet50) korištene su odgovarajuće metode predobrade iz biblioteke \texttt{tensorflow.keras.applications}, koje uključuju:
\begin{itemize}
    \item Središnjenje pikselskih vrijednosti
    \item Skaliranje prema statistici skupa na kojem je model treniran
\end{itemize}


\section{Modeli}

\subsection{Model 1: CNN}


\subsection{Model 2: MobileNetV2}



\subsection{Model 3: EfficientNetB0}



\subsection{Model 4: ResNet}


\newpage
\section{Usporedba modela}

\subsection{Korištene metrike}

Za kvantitativnu usporedbu korištene su sljedeće metrike:
\begin{itemize}
    \item \textbf{Točnost (Accuracy)} – ukupni broj točnih klasifikacija u odnosu na ukupan broj uzoraka.
    \item \textbf{Preciznost (Precision)} – omjer ispravno predviđenih pozitivnih uzoraka u odnosu na sve predviđene pozitivne uzorke.
    \item \textbf{Odziv (Recall)} – omjer ispravno predviđenih pozitivnih uzoraka u odnosu na sve stvarne pozitivne uzorke.
    \item \textbf{F1-mjera (F1 Score)} – harmonijska sredina preciznosti i odziva.
\end{itemize}

\subsection{Tablična usporedba rezultata}

\begin{table}[h]
\centering
\begin{tabular}{@{}lcccc@{}}
\toprule
Metrička mjera & CNN & ResNet & M3 & M4 \\
\midrule
Točnost (\%) & X1 & X2 & X3 & X4 \\
Preciznost (\%) & Y1 & Y2 & Y3 & Y4 \\
Odziv (\%) & Z1 & Z2 & Z3 & Z4 \\
F1-mjera (\%) & W1 & W2 & W3 & W4 \\
Vrijeme treniranja & T1 & T2 & T3 & T4 \\
Broj parametara & P1 & P2 & P3 & P4 \\
\bottomrule
\end{tabular}
\caption{Usporedba performansi modela prema različitim metrima}
\label{tab:rezultati_transponirano}
\end{table}



\subsection{Diskusija}

Nakon provođenja eksperimenta, može se primijetiti da pojedini modeli ostvaruju bolje rezultate ovisno o metrikama i resursima. Vrijeme treniranja i broj parametara značajno utječu na mogućnosti primjene u stvarnim okruženjima.

\section{Zaključak}

.

\begin{thebibliography}{25}

\bibitem{tensorflow}
TensorFlow documentation, \url{https://www.tensorflow.org/}

\bibitem{keras}
Keras API reference, \url{https://keras.io/api/}

\bibitem{mobilenet}
Sandler, M., et al. (2018). MobileNetV2: Inverted Residuals and Linear Bottlenecks. \textit{CVPR}.

\bibitem{efficientnet}
Tan, M., \& Le, Q. (2019). EfficientNet: Rethinking Model Scaling for Convolutional Neural Networks. \textit{ICML}.

\bibitem{resnet}
He, K., Zhang, X., Ren, S., \& Sun, J. (2016). Deep Residual Learning for Image Recognition. \textit{CVPR}.

\bibitem{asl_dataset}
DataMunge. 
Sign Language MNIST,
\url{https://www.kaggle.com/datasets/datamunge/sign-language-mnist}

\end{thebibliography}

\end{document}
